%
%
%%%%

%%%   Самарский государственный технический университет

%%%   Кафедра прикладной математики и информатики

%%%

%%%   Преамбула для статьи в журнал "Вестник Самарского государственного

%%%   технического университета. Серия: «Физико-математические науки»"

%%%   Ориентирована под MikTEx 2.4 for Win32

%%%

%%%   Хотя сам журнал собирается под GNU/Linux на дистрибутиве TeXLive



\documentclass[11pt,a4paper,twoside]{article}



\usepackage[cp1251]{inputenc}

\usepackage[T2A]{fontenc}

\usepackage[english,russian]{babel}

\usepackage{samgtubib}

\usepackage[dvips]{graphicx}

\usepackage[multiple]{footmisc}



\usepackage{samgtu}

\renewcommand{\VestnikYear}{2009} % Год издания Вестника

\renewcommand{\VestnikNumber}{2\,(19)} % Номер Вестника

\renewcommand{\VestnikMonth}{Ноябрь} % Месяц выхода Вестника

\renewcommand{\VestnikData}{01.11.09} % Дата подписания Вестника в печать

\renewcommand{\VestnikUPL}{26,64} % Усл. печ. л.

\renewcommand{\VestnikUIL}{25,73} % Уч.-изд. л.

\renewcommand{\VestnikRegNumber}{479/09} % Регистрационный номер

\renewcommand{\VestnikOrder}{994} % Номер заказа






%\usepackage{pst-barcode}
%\usepackage{auto-pst-pdf}
%
%\newcommand{\totop}[1]{\raisebox{6pt}[1pt][2pt]{#1}} 
%\newcommand{\tottop}[1]{\raisebox{12pt}[1pt][2pt]{#1}} 
%\newcommand{\totttop}[1]{\raisebox{18pt}[1pt][2pt]{#1}} 
%\newcommand{\tottttop}[1]{\raisebox{24pt}[1pt][2pt]{#1}} 
%\newcommand{\totttttop}[1]{\raisebox{30pt}[1pt][2pt]{#1}} 
%\newcommand{\tobot}[1]{\raisebox{-6pt}[1pt][2pt]{#1}} 
%\newcommand{\tobbot}[1]{\raisebox{-12pt}[1pt][2pt]{#1}} 
%\newcommand{\tobbbot}[1]{\raisebox{-18pt}[1pt][2pt]{#1}}
%
%%\samgtupubl % для генерации pdf, отдаваемого в типографию СамГТУ
%\pdfcrop % обрезка pdf для размещения на math-net.ru
%\draft % На корректуре
%%\nofilllastpage % Последняя русскоязычная страница заполнена не полностью
%\filllastpage
%%\briefcomm % Статья является кратким сообщением

\vsgtu{1321}

\FirstPage{44}
\LastPage{56}




%\begin{document}



\hfill \begin{pspicture}(1in,1in)
\psbarcode{http://dx.doi.org/10.14498/vsgtu1321}{}{qrcode}
\end{pspicture}
\vspace{-1in}


\udk{517.956.6}
\msc{35M12; 35A01, 35A02}

\rutitle{О разрешимости краевой задачи для~уравнения~смешанного типа с~сингулярным~коэффициентом}% Полный заголовок
{О разрешимости краевой задачи для уравнения смешанного типа \dots}

\ruauthor{М.~Х.~Рузиев}{Р\,у\,з\,и\,е\,в~М.~Х.}

\ruaffil{\rusimnuu.}
\enaffil{\engimnuu.}



\ruauthordetails{1}{\fullauthor{\rupid{42765}{Менглибай Холтожибаевич Рузиев}}\regalia{к.ф.-м.н.; \mem{mruziev@mail.ru}}\position{старший научный сотрудник.}}

\rukeywords{принцип экстремума, единственность решения, существование решения, интегральные уравнения, 
уравнение Винера"--~Хопфа, индекс уравнения, интеграл Фурье}

\ruabstract{Исследуется задача с условиями во внутренней характеристике и на частях  линии вырождения для уравнения смешанного типа  с сингулярным коэффициентом в неограниченной области. Единственность решения задачи доказывается с помощью принципа экстремума. При доказательстве существования решения задачи применяются теория сингулярных интегральных уравнений и интегральные уравнения Фредгольма.}

\entitle{On the Solvability of Boundary Value Problem for Mixed-type Equation with a Singular Coefficient}

\enauthor{M.~Kh.~Ruziev}{R\,u\,z\,i\,e\,v~M.~Kh.}


\enauthordetails{1}{\fullauthor{\enpid{42765}{Menglibay Kh. Ruziev}} \regalia{Cand. Phys. \& Math. Sci.; \mem{mruziev@mail.ru}} \position{Senior Scientific Researcher.} }

\enabstract{In this paper we study a problem with conditions on the inner characteristic and on some parts of the degeneration line for mixed type equation with singular coefficient in unbounded domain. We prove the uniqueness of solution of the mentioned problem with the help of the extremum principle. The proof of the existence of solution is based on the theory of singular integral equations and Fredholm integral equations.}

\enkeywords{principle of extremum, unique solvability, existence of solution, singular coefficient, integral equations, Wiener--Hopf equation, index equation, Fourier integral}

\date{25/III/2014}
\revisiondate{18/IV/2014}
\accepteddate{27/VIII/2014}

\makerutitle



\Section{Введение и постановка задачи}
Пусть $D={D^{+}}\cup D^{-}\cup I$ "--- область комплексной плоскости $z = x + iy$, 
где ${D^ + }$ "--- полуплоскость $y > 0$, $D^ {- }$ "--- конечная область полуплоскости $y < 0$, ограниченная характеристиками $AC$ и $BC$ уравнения
\begin{equation}
(\mathop{\rm sign}y){|y|^m}u_{xx}+u_{yy}+\frac{\beta_0}{y}u_y=0,
\label{ruz:1}
\end{equation}
исходящими из точек $A( - 1, 0)$, $B(1, 0)$, и отрезком $AB$ прямой $y = 0$; $I=\{(x,y) : -1<x<1, \; y=0\}$. 
В~уравнении \eqref{ruz:1} предполагается, что $m$, $\beta_0$ "--- некоторые действительные числа, удовлетворяющие
условиям $m > 0$, $- m/2 < \beta_0 < 1$.

Пусть $D_R^{+}$ "--- конечная область, отсекаемая от области $D^{+}$ дугой нормальной кривой $x^2+4y^{m + 2}/(m+2)^2=R^2$,
$- R \le x \le R$, $0 \le y \le \left((m+2)R/2\right)^{2/(m + 2)}$,
$A_R(-R, 0)$, $B_R(R, 0)$.

Введём обозначения: 
$$I_1=\{(x,y):-\infty  < x <- 1,\; y = 0\},\quad I_2=\{(x,y):1<x<\infty, \; y=0\},$$ 
$C_0(C_1)$ "--- точки пересечения характеристики $AC$ $(BC)$ с характеристикой, исходящей из
точки $E(c, 0)$, 
где $c \in I$ "--- произвольное фиксированное число, ${D_R} = D_R^{+}\cup D^-$, $D_R$ "--- подобласть неограниченной области $D$.

Пусть $p(x) \in {C^1}[ - 1,c]$ "--- диффеоморфизм, переводящий отрезок $[ - 1, c]$ в  отрезок $[c, 1]$, причём $p'(x) < 0$,
$p( - 1) = 1$, $p(c) = c$. 
В~качестве примера такой функции 
приведем линейную функцию $p(x) = \delta  - kx$, где
$k = (1 - c)/(1 + c)$, $\delta  = 2c/(1 + c)$,
$\delta  + k = 1$, $\delta  - kc = c$.

Краевая задача со смещением для уравнения смешанного типа в неограниченной области, эллиптическая часть которой "--- верхняя полуплоскость, исследована в работе \cite{ruz:bib:1}. Изучению краевой задачи в бесконечной полуполосе для обобщённого двуосесимметрического уравнения Гельмгольца посвящена работа \cite{ruz:bib:2}. Краевая задача для уравнения \eqref{ruz:1} в смешанной области, эллиптическая часть которой "--- полуполоса, решена в работе \cite{ruz:bib:3}.

Отметим, что в задаче Геллерстедта \cite{ruz:bib:4} значение искомой функции в гиперболической части смешанной области $D$
задается на характеристиках $E{C_0}$ и $E{C_1}$:
$${u \bigl|_{E{C_0}}} = {\psi _1}(x), \quad { u \bigr|_{E{C_1}}} = {\psi _2}(x).$$

В данной работе решается задача, где характеристика $E{C_1}$ освобождена от краевого условия, и это недостающее условие Геллерстедта заменено внутренне краевым условием локального смещения на отрезке
$AB$ линии вырождения $y = 0$.

\smallskip

\hypertarget{ruz:G}{}

\begin{newthm}{Задача $\Gamma $} Требуется найти в области $D$ функцию $u(x,y),$ удовлетворяющую следующим условиям\/{\rm :}

\begin{itemize}

\item[\rm 1)] функция $u(x,y)$ непрерывна в любой подобласти ${\bar D_R}$ неограниченной области $D;$

\item[\rm 2)] $u(x,y)$ принадлежит пространству ${C^2}({D^ + })$ и удовлетворяет уравнению \eqref{ruz:1} в этой области\/{\rm ;}

\item[\rm 3)] $u(x,y)$ является обобщенным решением класса ${R_1}[5]$  в области $D^ {- };$

\item[\rm 4)] выполняются равенства
\begin{equation}
\lim\limits_{R \to \infty } u(x,y) = 0, \quad y \ge 0, \quad {R^2} = {x^2} + 4{(m + 2)^{ - 2}}{y^{m + 2}};
\label{ruz:2}
\end{equation}

\item[\rm  5)]  $u(x,y)$ удовлетворяет краевым условиям\/{\rm :}
\begin{gather}
{\left. {u(x,y)} \right|_{y = 0}} = {\tau _i}(x), \quad \forall x \in {\bar I_i} \quad i = 1, 2,
\label{ruz:3}
\\
{\left. {u(x,y)} \right|_{E{C_0}}} = \psi (x), \quad (c - 1)/2 \le x \le c,
\label{ruz:4}
\\
u(p(x), 0)=\mu u(x, 0) + f(x), \quad  - 1 \le x \le c 
\label{ruz:5}
\end{gather}
и условию сопряжения
\begin{equation}
\lim\limits_{y \to+0} {y^{\beta_0}}{u_y}=\lim\limits_{y\to-0}{(-y)^{\beta_0}}{u_y}, \quad x \in I\backslash \{ c\},
\label{ruz:6}
\end{equation}
причём эти пределы при $x = \pm 1,$ $x = c$ могут иметь особенности
порядка ниже $1 - 2\beta,$ где $\beta  = (m + 2{\beta _0})/(2(m + 2)),$
$f(x),$ $\psi (x),$ ${\tau _i}(x),$ $i = 1, 2$ "--- заданные функции,
причём $f(x) \in C[ - 1,c] \cap {C^{(1,{\alpha _0})}}( - 1,c),$
$f(c) = 0,$  $f( - 1) = 0,$ $\psi (x) \in C[(c - 1)/2,c] \cap {C^{(1,{\delta _0})}}((c - 1)/2,c),$
$\psi (c) = 0,$ $\mu$ "--- ${\rm const},$
функции ${\tau _i}(x)$ в окрестности точек $x=-1$,  $x = 1$ представимы в виде
 $\tau _i(x) = (1 - x^2){\tilde \tau _i}(x)$ и они удовлетворяют условию Гельдера на любых интервалах 
  $( - N, - 1),$ $(1,N),$ $N > 1$ и  для  достаточно больших $|x|$
  удовлетворяют неравенству $|\tau_i(x)|\leq M{|x|^{-\delta}},$
  где $\delta,$~$M$ "--- положительные постоянные{\rm .}
  
\end{itemize}
\end{newthm}

Отметим, что условие \eqref{ruz:5} является внутренне краевым условием локального
смещения на отрезке линии параболического вырождения [\citen{ruz:bib:6}--\citen{ruz:bib:8}].


\smallskip

\Section{Единственность решения задачи \hyperlink{ruz:G}{$\Gamma $}}

\begin{theorem}[1]Пусть выполнены условия  $Х{\tau _i}(x) \equiv 0,$
$i = 1, 2,$  $\psi (x) \equiv 0,$  ${f(x) \equiv 0,}$  $0 < \mu  < 1.$
Тогда задача \hyperlink{ruz:G}{$\Gamma $}  имеет  лишь тривиальное решение{\rm .}
\end{theorem}

\smallskip

\begin{proof} Решение видоизмененной задачи Коши для
уравнения \eqref{ruz:1} в области $D^{-}$, удовлетворяющее начальным
данным  
$$\lim\limits_{y\to-0}u(x,y)=\tau(x),\quad x\in\bar I;\qquad \lim\limits_{y\to-0}{(-y)^{\beta_0}}u_{y}=\nu(x),\quad  x\in I,$$ даётся формулой Дарбу
\cite[с.~34]{ruz:bib:9}:
\begin{multline}
u(x,y)=\gamma_1\int_{-1}^1{\tau\Bigl[x+\frac{2t}{m + 2}(-y)^{(m+2)/2}\Bigr](1+t)^{\beta-1}(1-t)^{\beta-1}}dt+
\\
+\gamma_2(-y)^{1-\beta_0}\int_{-1}^1{\nu\Bigl[x+\frac{2t}{m+2}(-y)^{(m + 2)/2}\Bigr](1+t)^{-\beta}(1-t)^{-\beta}dt},
\label{ruz:7}
\end{multline}
где
$\gamma_1=\Gamma(2\beta){2^{1-2\beta}}/\Gamma^2(\beta)$,
$\gamma_2=-\Gamma(2-2\beta){2^{2\beta-1}}/(1-\beta_0)\Gamma^2(1-\beta)$,
$\Gamma(z)$ "--- гамма"=функция \cite{ruz:bib:4}.

В силу формулы \eqref{ruz:7}  из краевого условия \eqref{ruz:4} после несложных вычислений получим
\begin{equation}
\nu(X)=\gamma D_{X,c}^{1-2\beta}\tau(X)+\Psi(X), \quad X\in(-1,c),
\label{ruz:8}
\end{equation}
где
$$\Psi(X)=\frac{(c-X)^\beta D_{X,c}^{1-\beta}\psi((X+c)/2)}{\gamma_2(m+2)/2)^{1-2\beta}\Gamma(1-\beta)},\quad X=2x-c,$$ $$\gamma=\frac{2\Gamma(1-\beta)\Gamma(2\beta)((m+2)/4)^{2\beta}}{\Gamma(\beta)\Gamma(1-2\beta)},$$  $D_{X,c}^l$ "--- оператор дробного дифференцирования в смысле Римана"--~Лиувилля~\cite{ruz:bib:5}.

Равенство \eqref{ruz:8} является первым функциональным соотношением между
неизвестными функциями  $\tau (x)$ и $\nu (x)$, принесённым на
интервал $(-1, c)$  оси $y=0$ из  гиперболической части $D^{ - }$
смешанной области $D$.

Теперь докажем, что если
$$\tau_i(x)\equiv0, \quad i=1, 2, \quad \psi(x)\equiv0, \quad f(x)\equiv0, \quad 0<\mu<1,$$
то решение задачи \hyperlink{ruz:G}{$\Gamma $} в области $D^{+}\cup I_1\cup\bar I\cup I_2$ в силу \eqref{ruz:2} тождественно равно нулю.

Пусть $(x_0, y_0)$ "--- точка положительного максимума функции $u(x,y)$ в области $\bar D_R^{+}$.

В силу \eqref{ruz:2}  $\forall\varepsilon>0$ существует такое $R_0=R_0(\varepsilon)$, что при $R>R_0(\varepsilon)$
\begin{equation}
|u(x,y)|<\varepsilon, \quad (x,y)\in A_RB_R.
\label{ruz:9}
\end{equation}
В силу обозначения $u(x, 0)=\tau(x)$, $x\in\bar I$ условие \eqref{ruz:5} перепишем в виде
\begin{equation}
\tau(p(x))=\mu\tau(x)+f(x), \quad x\in[-1,c].
\label{ruz:10}
\end{equation}

Отсюда при $x=c$, (где $f(x)\equiv0$), имеем $\tau(p(c))=\mu\tau(c)$.
Тогда в силу равенства $p(c)=c$ следует, что $\tau(c)(1-\mu)=0$,
т.~е. $\tau(c)=0$. В силу принципа Хопфа \cite[с.~25]{ruz:bib:10}
функция $u(x,y)$ своего положительного максимума и~отрицательного минимума во внутренних точках
области $\bar D_R^ + $ не достигает. В~силу  $0<\mu<1$
из \eqref{ruz:10} (где $f(x)\equiv0$) следует, что их также нет и в~интервале $(c, 1)$ оси $y = 0$. 
Допустим, что искомая функция своего положительного максимума и отрицательного минимума
достигает в точках интервала $(-1, c)$ оси $y = 0$.

Пусть $(x_0, 0)$ (где $x_0\in(-1, c)$) "---  точка положительного максимума (отрицательного минимума)
функции $u(x, 0)=\tau(x)$. 
Тогда в этой точке в случае положительного максимума (отрицательного минимума) \cite[c.~74]{ruz:bib:9}
\begin{equation}
\nu(x_0)<0 \quad (\nu(x_0)>0).
\label{ruz:11}
\end{equation}

Хорошо известно, что в точке положительного максимума
(отрицательного минимума) функции $\tau (x)$ для операторов
дробного дифференцирования имеет место неравенство
$D_{x_0,c}^{1-2\beta}\tau(x)>0$ $(D_{x_0,c}^{1-2\beta}\tau(x)<0)$.
Тогда в силу \eqref{ruz:8} (где $\Psi(x)\equiv0$) 
\begin{equation}
\nu(x_0)=\gamma D_{x_0,c}^{1-2\beta}\tau(x)>0 \quad
(\nu(x_0)=\gamma D_{x_0,c}^{1-2\beta}<0).
\label{ruz:12}
\end{equation}

Неравенства \eqref{ruz:11} и \eqref{ruz:12} противоречат условию сопряжения \eqref{ruz:6},
отсюда следует, что ${x_0}\notin(-1, c)$.

Следовательно, точки положительного максимума (отрицательного минимума)
функции $u(x,y)$ нет на интервале $AB$. Пусть $R>{R_0}$.
Из принципа Хопфа и предыдущих рассуждений получаем, что
$(x_0,y_0)\in A_RB_R$  и в силу \eqref{ruz:9} "--- $|u(x_0,y_0)|<\varepsilon$.
Следовательно, $|u(x,y)|<\varepsilon$
$\forall (x,y)\in\bar D_R^{+}$. Отсюда в силу произвольности
$\varepsilon$  при $R\to+\infty$ заключаем, что $u(x,y)\equiv0$
в области $D^{+}\cup I_1\cup\bar I\cup I_2$. Тогда
\begin{equation}
\lim\limits_{y\to+0}u(x,y)=0, \quad x\in\bar I; \quad \lim\limits_{y\to+0}y^{\beta_0}u_y=0, \quad x\in I.
\label{ruz:13}
\end{equation}

С учётом \eqref{ruz:13} в силу непрерывности решения в области
$\bar D_R^{+}$ и условия сопряжения \eqref{ruz:6} восстанавливая искомую
функцию $u(x,y)$ в области $D^-$ как решение  видоизмененной задачи
Коши с однородными данными, получим $u(x,y)\equiv0$ в
области $\bar D^-$. \end{proof}

\smallskip

\Section{Существование решения задачи \hyperlink{ruz:G}{$\Gamma $}}

\begin{theorem}[2]Пусть выполнены условия $p(x){=}\delta-kx,$ $\mu k^{{1}/{2}-3\alpha}<1,$
${0<\mu<1},$ $\beta_0>({1-m})/{3},$
где  $\delta=2c/(1+c),$ $k=(1-c)/(1+c),$ $\alpha=(1-2\beta)/4.$ Тогда решение задачи \hyperlink{ruz:G}{$\Gamma $} существует{\rm .}
\end{theorem}

\begin{proof} Решение задачи Дирихле, удовлетворяющее
условиям \eqref{ruz:3} и $u(x, 0)=\tau (x)$, $x\in\bar I$, представимо в виде
\begin{equation}
u(x,y)=k_2(1-\beta_0)y^{1-\beta_0}\int_{-1}^1{\tau(t)(r_0^2)^{\beta-1}}dt+F_1(x,y),
\label{ruz:14}
\end{equation}
где 
\begin{gather*}
r_0^2=(x-t)^2+\frac{4}{(m+2)^2}y^{m+2},
\\
F_1(x,y)=k_2(1-\beta_0)y^{1-\beta_0}\left[\int_{-\infty}^{-1}{\tau_1(t)(r_0^2)^{\beta-1}}dt+\int_1^\infty{\tau_2(t)(r_0^2)^{\beta-1}dt}\right],
\\
k_2=\frac{1}{4\pi}\left(\frac{4}{m+2}\right)^{2-2\beta}\frac{\Gamma^2(1-\beta)}{\Gamma(2-2\beta)}.
\end{gather*}

Дифференцируя \eqref{ruz:14} по $y$ и учитывая равенство
\begin{multline*}
\frac{\partial}{\partial y}\left\{y^{1-\beta_0}\Bigl[(x-t)^2+\frac{4}{(m+2)^2}y^{m+2}\Bigr]^{\beta-1}\right\}=
\\
=\frac{m+2}{2}y^{-\beta_0}\frac{\partial}{\partial t}\left\{(x-t)\Bigl[(x-t)^2+\frac{4}{(m+2)^2}y^{m+2}\Bigr]^{\beta-1}\right\},
\end{multline*}
получим
\begin{multline}
\frac{\partial u}{\partial y}=k_2(1-\beta_0)\frac{m+2}{2}y^{-\beta_0}\times
\\
\times\int_{-1}^1{\tau(t)\frac{\partial}{\partial t}\left\{(x-t)
\Bigl[(x-t)^2+\frac{4}{(m+2)^2}y^{m+2}\Bigr]^{\beta-1}\right\}dt}+
\frac{\partial F_1(x,y)}{\partial y}.
\label{ruz:15}
\end{multline}

В интеграле правой части  равенства \eqref{ruz:15}, выполнив операцию интегрирования по частям,
с учётом $\tau(-1)=0$,  $\tau(1)=0$ после несложных вычислений имеем 
\begin{multline}
\frac{\partial u}{\partial y}=-k_2(1-\beta_0)\frac{m+2}{2}y^{-\beta_0}\int_{-1}^1{\tau'(t)(x-t)\Bigl[(x-t)^2+\frac{4}{(m+2)^2}y^{m+2}\Bigr]^{\beta-1}}dt+
\\
+\frac{\partial F_1(x,y)}{\partial y}.
\label{ruz:16}
\end{multline}

Умножая обе части равенства \eqref{ruz:16}  на  $y^{\beta_0}$ и затем  переходя к пределу при $y\to+0$, получим
\begin{equation}
\nu(x)=k_2(1-\beta_0)\frac{m+2}{2}\int_{-1}^1{\frac{(x-t)\tau'(t)dt}{|x-t|^{2-2\beta}}}+\Phi(x), \quad x\in(-1, 1),
\label{ruz:17}
\end{equation}
где
$$\Phi(x)=\lim\limits_{y\to+0}y^{\beta_0}\frac{\partial F_1(x,y)}{\partial y}=k_2(1-\beta_0)^2\left(\int_{-\infty}^{-1}{\frac{\tau_1(t)dt}{(x-t)^{2-2\beta}}}+\int_1^\infty{\frac{\tau_2(t)dt}{(t-x)^{2-2\beta}}}\right).$$

Это есть второе функциональное соотношение между неизвестными
функциями $\nu(x)$ и $\tau(x)$, принесёнными на интервал
$I$ оси $y=0$ из  верхней полуплоскости.

Заметим, что соотношение \eqref{ruz:1} справедливо для всего промежутка $I$.

Далее, разбивая промежуток интегрирования $(-1, 1)$  на промежутки \linebreak $(-1, c)$ и $(c, 1)$, а затем  в интегралах с пределом
$(c, 1)$ сделав замену переменного интегрирования $t=p(s)=\delta-ks$,
учитывая равенство \eqref{ruz:10}, соотношение \eqref{ruz:17} приведём к виду
\begin{multline}
\nu(x)=-k_2(1-\beta_0)\frac{m+2}{2}\left[\int_{-1}^x{\frac{\tau'(t)dt}{(x-t)^{1-2\beta}}}-\int_x^c{\frac{\tau'(t)dt}{(t-x)^{1-2\beta}}}+\right.
\\
+\left.\mu\int_{-1}^c{\frac{\tau'(s)ds}{(p(s)-x)^{1-2\beta}}}+\int_{-1}^c{\frac{f'(s)ds}{(p(s)-x)^{1-2\beta}}}\right]+\Phi(x), \quad x\in( - 1, c).
\label{ruz:18}
\end{multline}
В силу \eqref{ruz:6}, исключая функцию $\nu(x)$ из \eqref{ruz:8} и \eqref{ruz:18}, получим
\begin{multline}
\frac{-2\gamma}{k_2(1-\beta_0)(m+2)}D_{x,c}^{1-2\beta}\tau(x)+F_0(x)=
\\
=\mu\int_{-1}^c{\frac{\tau'(s)ds}{(p(s)-x)^{1-2\beta}}}+\int_{-1}^x{\frac{\tau'(t)dt}{(x-t)^{1-2\beta}}}-\int_x^c{\frac{\tau'(t)dt}{(t-x)^{1-2\beta}}}, \quad x\in(-1, c),
\label{ruz:19}
\end{multline}
где
$$F_0(x)=-\frac{2(\Psi(x)-\Phi(x))}{k_2(1-\beta_0)(m+2)}-\int_{-1}^c{\frac{f'(s)ds}{(p(s)-x)^{1-2\beta}}}.$$

Применив  оператор $\Gamma(1-2\beta)D_{x,c}^{2\beta-1}$ к обеим частям равенства \eqref{ruz:19} и учитывая, 
что $D_{x,c}^{2\beta-1}D_{x,c}^{1-2\beta}\tau(x)=\tau(x)$, имеем
\begin{multline}
\frac{-2\gamma}{k_2(1-\beta_0)(m+2)}\Gamma(1-2\beta)\tau(x)+\Gamma(1-2\beta)D_{x,c}^{2\beta-1}F_0(x)=
\\
=\Gamma(1-2\beta)D_{x,c}^{2\beta-1}\biggr[\mu\int_{-1}^c{\frac{\tau'(s)ds}{(p(s)-x)^{1-2\beta}}}+\int_{-1}^x{\frac{\tau'(t)dt}{(x-t)^{1-2\beta}}}- \\ -
\int_x^c{\frac{\tau'(t)dt}{(t-x)^{1-2\beta}}}\biggl], 
\quad 
x\in[-1, c].
\label{ruz:20}
\end{multline}
Далее нетрудно убедиться в том, что
\begin{multline}
\Gamma(1-2\beta)D_{x,c}^{2\beta-1}\int_{-1}^x{\frac{\tau'(t)dt}{(x-t)^{1-2\beta}}} = \\
=-\pi\ctg(2\beta\pi)\tau(x)+\int_{-1}^c{\Bigl(\frac{c-x}{c-t}\Bigr)^{1-2\beta}\frac{\tau(t)dt}{t - x}},
\label{ruz:21}
\end{multline}
\begin{equation}
\Gamma(1-2\beta)D_{x,c}^{2\beta-1}\int_x^c{\frac{\tau'(t)dt}{(t-x)^{1-2\beta}}}=-\frac{\pi}{\sin(2\beta\pi)}\tau(x),
\label{ruz:23}
\end{equation}
\begin{multline}
\Gamma(1-2\beta)D_{x,c}^{2\beta-1}\mu\int_{-1}^c{\frac{\tau'(s)ds}{(p(s)-x)^{1-2\beta}}} = \\
=\mu\int_{-1}^c{\Bigl(\frac{c-x}{p(s)-c}\Bigr)^{1-2\beta}\frac{\tau(s)p'(s)ds}{p(s)-x}}.
\label{ruz:23}
\end{multline}

\pagebreak

Подставляя \eqref{ruz:21}--\eqref{ruz:23} в \eqref{ruz:20}, после несложных вычислений получим
сингулярное интегральное уравнение относительно $\tau(x)$:
\begin{multline}
\tau(x)+\lambda\int_{-1}^c{\Bigl(\frac{c-x}{c-t}\Bigr)^{1-2\beta}\frac{\tau(t)dt}{t-x}}=
\\
=-\lambda\mu\int_{-1}^c{\Bigl(\frac{c-x}{p(s)-c}\Bigr)^{1-2\beta}\frac{\tau(s)p'(s)ds}{p(s)-x}}+F_1(x), \quad x\in[-1, c],
\label{ruz:24}
\end{multline}
где 
\begin{gather*}
F_1(x)=\lambda\Gamma(1-2\beta)D_{x,c}^{2\beta-1}F_0(x),
\\
F_1(x)\in C[-1,c]\cap C^{0,\bar\gamma}(-1,c), \quad \bar\gamma>1-\beta, \quad \lambda=\frac{\cos(\beta\pi)}{\pi(1+\sin(\beta\pi))}.
\end{gather*}

Интегральный оператор правой части равенства \eqref{ruz:24} не является регулярным,
так как подынтегральное выражение при $x=c$, $s=c$ имеет изолированную
особенность первого порядка, поэтому это слагаемое в \eqref{ruz:24} выделено отдельно.

Временно считая правую часть уравнения \eqref{ruz:24} известной функцией,
перепишем его в виде
\begin{equation}
\tau(x)+\lambda\int_{-1}^c{\Bigl(\frac{c-x}{c-t}\Bigr)^{1-2\beta}\frac{\tau(t)dt}{t-x}}=g_0(x), \quad x\in[-1, c],
\label{ruz:25}
\end{equation}
где
\begin{equation}
g_0(x)=-\lambda\mu\int_{-1}^c{\Bigl(\frac{c-x}{p(s)-c}\Bigr)^{1-2\beta}\frac{\tau(s)p'(s)ds}{p(s)-x}}+F_1(x).
\label{ruz:26}
\end{equation}

Полагая $(c-x)^{2\beta-1}\tau(x)=\rho(x)$, $(c-x)^{2\beta-1}g_0(x)=g_1(x)$,
уравнение \eqref{ruz:25} запишем в виде
\begin{equation}
\rho(x)+\lambda\int_{-1}^c{\frac{\rho(t)dt}{t-x}}=g_1(x).
\label{ruz:27}
\end{equation}

Решение уравнения \eqref{ruz:27} будем искать в классе функций, удовлетворяющих
условию Гёльдера на  $(-1, c)$ и ограниченных при $x=-1$,
а при $x=c$ могущих обращаться в бесконечность порядка меньше $1-2\beta$. 
В~этом классе индекс уравнения \eqref{ruz:27}
равен нулю.  Решение уравнения \eqref{ruz:27} находится в явном виде методом Карлемана"--~Векуа \cite{ruz:bib:11}:
$$\rho(x)=\frac{1+\sin(\beta\pi)}{2}g_1(x)-
\frac{\cos(\beta\pi)}{2\pi}\Bigl(\frac{1+x}{c-x}\Bigr)^{(1-2\beta)/4}\int_{-1}^c{\Bigl(\frac{c-t}{1+t}\Bigr)^{(1-2\beta)/4}\frac{g_1(t)dt}{t-x}}.$$

Отсюда, возвращаясь к прежним функциям, получим
\begin{equation}
\tau(x)=\cos^2(\pi\alpha)g_0(x)-\frac{\sin(2\pi\alpha)}{2\pi}\int_{-1}^c{\frac{(1+x)^\alpha(c-x)^{3\alpha}}{(1+t)^\alpha(c-t)^{3\alpha}}\frac{g_0(t)dt}{t-x}},
\label{ruz:28}
\end{equation}
где $\alpha=(1-2\beta)/4$.

Теперь, подставляя \eqref{ruz:26} в \eqref{ruz:28}, после некоторых преобразований имеем
\begin{multline}
\tau(x)=-\lambda\mu\cos^2(\pi\alpha)\int_{-1}^c{\Bigl(\frac{c-x}{p(s)-c}\Bigr)^{4\alpha}\frac{\tau(s)p'(s)ds}{p(s)-x}}+
\\
+\lambda\mu\frac{\sin(2\pi\alpha)}{2\pi}(1+x)^\alpha(c-x)^{3\alpha}\int_{-1}^c{\frac{\tau(s)p'(s)ds}{(p(s)-c)^{4\alpha}}}\times
\\
\times\int_{-1}^c{\Bigl(\frac{c-t}{1+t}\Bigr)^\alpha\frac{dt}{(p(s)-t)(t-x)}}+F_2(x), \quad x\in[-1, c],
\label{ruz:29}
\end{multline}
где
$$
F_2(x)=\cos^2(\pi\alpha)F_1(x)-\frac{\sin(2\pi\alpha)}{2\pi}\int_{-1}^c{\Bigl(\frac{1+x}{1+t}\Bigr)^\alpha
\Bigl(\frac{c-x}{c-t}\Bigr)^{3\alpha}\frac{F_1(t)dt}{t-x}}.
$$

В силу $p(x)=\delta-kx$, где $k=(1-c)/(1+c)$, $\delta=2c/(1+c)$, уравнение \eqref{ruz:29}  запишем в виде
\begin{multline}
\tau(x)=\lambda\mu k^{1-4\alpha}\cos^2(\pi\alpha)\int_{-1}^c{\Bigl(\frac{c-x}{c-s}\Bigr)^{4\alpha}\frac{\tau(s)ds}{\delta-ks-x}}-
\\
-\lambda\mu k^{1-4\alpha}\frac{\sin(2\pi\alpha)}{2\pi}(1+x)^\alpha(c-x)^{3\alpha}\times
\\
\times\int_{-1}^c{\frac{\tau(s)ds}{(c-s)^{4\alpha}}}\int_{-1}^c{\Bigl(\frac{c-t}{1+t}\Bigr)^\alpha\frac{dt}{(t-x)(\delta-ks-t)}}+F_2(x), \quad x\in[-1,c].
\label{ruz:30}
\end{multline}

В \eqref{ruz:30} вычислим внутренний интеграл:
$$A(x,s)=\int_{-1}^c{\Bigl(\frac{c-t}{1+t}\Bigr)^\alpha\frac{dt}{(t-x)(\delta-ks-t)}}.$$
Разлагая рациональный множитель подынтегрального выражения на простые дроби, используя формулы гипергеометрической функции  и выполнив несложные вычисления, имеем
\begin{multline}
A(x,s)=\frac{1}{\delta-ks-x}\Bigl[\pi\ctg(\pi\alpha)\frac{(c-x)^\alpha}{(1+x)^\alpha}+\Gamma(-\alpha)\Gamma(1+\alpha)+
\\
+\frac{1+c}{1+\delta-ks}\Gamma(1+\alpha)\Gamma(1-\alpha) F\Bigl(1-\alpha, 1, 2;\frac{1+c}{1+\delta-ks}\Bigr)\Bigr],
\label{ruz:31}
\end{multline}
где
$F(a,b,c;z)$ "--- гипергеометрическая функция Гаусса \cite{ruz:bib:5}.

Подставляя \eqref{ruz:31} в \eqref{ruz:30}, после несложных вычислений получим следующее интегральное уравнение:
\begin{equation}
\tau(x)=\lambda\int_{-1}^c{\frac{{\rm K}(x,s)\tau(s)ds}{\delta-ks-x}}+F_2(x), \quad x\in[-1, c],
\label{ruz:32}
\end{equation}
где
\begin{multline*}
K(x,s)=-\mu k^{1-4\alpha}\frac{\sin(2\pi\alpha)}{2\pi}\frac{(1+x)^\alpha(c-x)^{3\alpha}}{(c-s)^{4\alpha}}\times 
\\
\times\left[\Gamma(-\alpha)\Gamma(1+\alpha)+\frac{1+c}{1+\delta-ks}\Gamma(1+\alpha)\Gamma(1-\alpha)F\left(1-\alpha, 1, 2;\frac{1+c}{1+\delta-ks}\right)\right].
\end{multline*}
\pagebreak

Применив  формулу  Больца  \cite[с.~11]{ruz:bib:5} для гипергеометрической функции
$$F\Bigl(1-\alpha, 1, 2;\frac{1+c}{1+\delta-ks}\Bigr),$$ 
в силу формул 
$$F(a,b,b;z)=(1-z)^{-a}~\text{\cite[с.~13]{ruz:bib:5},}$$
$$
\Gamma(1+z)=z\Gamma(z),\; \; \Gamma(z)\Gamma(1-z)=\frac{\pi}{\sin(\pi z)}~\text{\cite[с.~5]{ruz:bib:5},}$$ функцию ${\rm K}(x,s)$ запишем в виде
\begin{equation}
{\rm K}(x,s)=\mu k^{1-3\alpha}\cos(\pi\alpha)
\Bigl(\frac{1+x}{1+\delta-ks}\Bigr)^\alpha\Bigl(\frac{c-x}{c-s}\Bigr)^{3\alpha}.
\label{ruz:33}
\end{equation}
Подставив \eqref{ruz:33} в равенство \eqref{ruz:32}, получим
\begin{multline}
\tau(x)=\lambda\mu k^{1-3\alpha}\cos(\pi\alpha)\int_{-1}^c{\Bigl(\frac{1+x}{1+\delta-ks}\Bigr)^\alpha
\Bigl(\frac{c-x}{c-s}\Bigr)^{3\alpha}
\frac{\tau(s)ds}{\delta-ks-x}} +\\ +F_2(x),\quad  x\in[-1,c].
\label{ruz:34}
\end{multline}
Выделив в уравнении \eqref{ruz:34}  характеристическую часть, преобразуем его к виду
\begin{multline}
\tau(x)=\lambda\mu k^{1-3\alpha}\cos(\pi\alpha)\int_{-1}^c 
\Bigl(\frac{c-x}{c-s}\Bigr)^{3\alpha}\frac{\tau(s)ds}{\delta-ks-x}+
\\
+R_1[\tau(x)]+F_2(x),\quad x\in[-1, c],
\label{ruz:35}
\end{multline}
где
$$R_1[\tau(x)]=\lambda\mu k^{1-3\alpha}\cos(\pi\alpha)\int_{-1}^c
\Bigl(\frac{c-x}{c-s}\Bigr)^{3\alpha}
\frac{\tau(s)}{\delta-ks-x}\left[\Bigl(\frac{1+x}{1+\delta-ks}\Bigr)^\alpha-1\right]ds$$
"--- регулярный оператор. Уравнение \eqref{ruz:35} запишем в виде
\begin{multline}
\tau(x)=\lambda\mu k^{1-3\alpha}\cos(\pi\alpha)\int_{-1}^c
\Bigl(\frac{c-x}{c-s}\Bigr)^{3\alpha}\frac{\tau(s)ds}{(c-s)
\bigl(k+({c-x})/({c-s})\bigr)}+
\\
+R_1[\tau(x)]+F_2(x),\quad x\in[-1, c].
\label{ruz:36}
\end{multline}
Осуществляя в равенстве \eqref{ruz:36} замену переменных  
$$x=c-(1+c){e^{-\xi}},\quad s=c-(1+c){e^{-t}}$$ и обозначая
$$\rho(\xi)=\tau(c-(1+c){e^{-\xi}}){e^{(3\alpha-1/2)\xi}},$$ 
приведём его к виду
\begin{multline}
\rho(\xi)=\lambda\mu k^{1-3\alpha}\cos(\pi\alpha)\int_0^\infty\frac{\rho(t)dt}{ke^{({\xi-t})/{2}}+e^{-({\xi-t})/{2}}}+
\\
+R_2[\rho(\xi)]+F_3(\xi),\quad \xi\in(0, \infty),
\label{ruz:37}
\end{multline} \pagebreak

\noindent
где
$$R_2[\rho(\xi)]=R_1[\tau(c-(1+c)e^{-\xi})]e^{(3\alpha-1/2)\xi},$$
$$F_3(\xi)=F_2(c-(1+c){e^{-\xi}}){e^{(3\alpha-1/2)\xi}}.$$
Заметим, что в силу условия $3\beta_0>{1-m}$ имеет место неравенство $6\alpha-{1}<0$.

Введём обозначение  
$${\rm N}(\zeta)=\frac{\lambda\mu k^{1-3\alpha}\cos(\pi\alpha)}{ke^{\zeta/2}+e^{-\zeta/2}}.$$
Тогда  уравнение \eqref{ruz:37} запишем в виде
\begin{equation}
\rho(\xi)=\int_0^\infty{{\rm N}(\xi-t)\rho(t)dt}+R_2[\rho(\xi)]+F_3(\xi),\quad \xi\in(0,\infty ). 
\label{ruz:38}
\end{equation}

Уравнение \eqref{ruz:38} является  интегральным  уравнением Винера"--~Хопфа \cite[с.~55]{ruz:bib:12}
и с помощью преобразования Фурье оно приводится к краевой задаче Римана, т.~е. решается в квадратурах.
Функции ${\rm N}(\xi)$, $F_3(\xi)$ имеют показательный порядок убывания на бесконечности, причём
$${\rm N}'(\xi)\in C(0, \infty),\quad  F_3(\xi)\in {{\rm H}^{{\alpha _1}}}(0,\infty ).$$
Следовательно, ${\rm N}(\xi )$, ${F_3}(\xi ) \in {L_2} \cap {{\rm H}^{{\alpha _1}}}$,
а решение уравнения \eqref{ruz:38} ищется в классе $\{0\}$ \cite[с.~12]{ruz:bib:12}.
Теоремы Фредгольма для интегральных уравнений типа свёртки справедливы лишь в одном частном случае "--- когда индекс
этих уравнений равен нулю.
Индексом уравнения \eqref{ruz:38} будет индекс выражения \mbox{$1-\hat N(\xi)$} с обратным знаком,
где
\begin{equation}
\hat N (\xi)=\int_{-\infty}^\infty{e^{i\xi t}{\rm N}(t)dt}=\lambda\mu k^{1-3\alpha}\cos(\pi\alpha)\int_{-\infty}^\infty{\frac{e^{i\xi t}dt}{ke^{t/2}+e^{-t/2}}}.
\label{ruz:39}
\end{equation}
Вычислив интеграл Фурье, с помощью теории вычетов \cite[с.~198]{ruz:bib:9} найдём
\begin{equation}
\int_{-\infty}^\infty{\frac{e^{i\xi t}dt}{ke^{t/2}+e^{-t/2}}}=\frac{\pi e^{-i\xi\ln k}}{\sqrt k \ch(\pi\xi)}.
\label{ruz:40}
\end{equation}
Подставляя \eqref{ruz:40} в  равенство \eqref{ruz:39}, используя
$$\pi\lambda\cos(\pi\alpha)=\pi\frac{\cos(\beta\pi)}{\pi(1+\sin(\beta\pi))}\cos(\alpha\pi)=\sin(\alpha\pi),$$
имеем
$$ \hat N (\xi)=\mu k^{{1}/{2}-3\alpha}\sin(\pi\alpha)\frac{e^{-i\xi\ln k}}{\ch(\pi\xi)}.$$
В силу условия
$$\mu k^{{1}/{2}-3\alpha}<1$$ 
и так как
\begin{multline*}
\mathop{\rm Re}(\hat N(\xi))=\mathop{\rm Re}\left(\mu k^{{1}/{2}-3\alpha}\sin(\pi\alpha)\frac{e^{-i\xi\ln k}}{\ch(\pi\xi)}\right)=
\\
=\mu k^{{1}/{2}-3\alpha}\sin(\pi\alpha)\frac{\cos(\xi\ln k)}{\ch(\pi\xi)}<\mu k^{{1}/{2}-3\alpha}<1,
\end{multline*}
 $\mathop{\rm Re}(1-\hat N(\xi))>0$.
Следовательно, индекс уравнения \eqref{ruz:38} $$\chi=-\mathop{\rm Ind}(1-\hat N(\xi))=0,$$
 т.~е. изменение аргумента выражения $1-\hat N (\xi)$ на действительной оси, выраженное в полных
 оборотах, равно нулю \cite[с.~56]{ruz:bib:12}. 
 Следовательно, уравнение \eqref{ruz:38} редуцируется к интегральному уравнению Фредгольма
  второго рода, однозначная разрешимость которого следует из единственности решения задачи~\hyperlink{ruz:G}{$\Gamma $}.~\end{proof}



\thanks{{\bf ORCID}

{\bf Menglibay Ruziev:} \url{http://orcid.org/0000-0002-1097-0137}
}

\RusRef


\begin{thebibliography}{99}


\RBibitem{ruz:bib:1} 
\by Репин~О.~А., Кумыкова~С.~К.
\paper Об одной краевой задаче со смещением для уравнения смешанного типа в неограниченной области
\jour Диффер. уравн.
\yr 2012
\vol 48
\issue 8 
\pages 1140--1149
\zmath{https://zbmath.org/?q=an:1260.35080}
%Repin, O.A.; Kumykova, S.K.
%On a boundary value problem with shift for an equation of mixed type in an unbounded domain. (English. Russian original) Zbl 1260.35080
%Differ. Equ. 48, No. 8, 1127-1136 (2012); translation from Differ. Uravn. 48, No. 8, 1140-1149 (2012).


\RBibitem{ruz:bib:2} 
\by Абашкин~А.~А.
\yr 2012
\issue 1(26) 
\pages 39--45
\paper Об одной задаче для обобщённого двуосесимметрического уравнения Гельмгольца в~бесконечной полуполосе
\jour Вестн. Сам. гос. техн. ун-та. Сер. Физ.-мат. науки
\mathnet{http://mi.mathnet.ru/vsgtu1023}
\crossref{10.14498/vsgtu1023}


\RBibitem{ruz:bib:3} 
\by Рузиев~М.~Х. 
\yr 2009 
\issue 1(18) 
\pages 33--40
\paper Краевая задача для уравнения смешанного типа с~сингулярным коэффициентом в области, эллиптическая часть которой~--- полуполоса
\jour Вестн. Сам. гос. техн. ун-та. Сер. Физ.-мат. науки
\mathnet{http://mi.mathnet.ru/vsgtu645}
\crossref{10.14498/vsgtu645}



\Bibitem{ruz:bib:4} 
\by Gellerstedt~S. 
\paper Quelques probl\`emes mixtes pour l'\'equation ${y^m}{z_{xx}} + {z_{yy}} = 0$ 
\jour Ark. Mat. Astron. Fys.
\yr 1937
\vol 26A
\issue 3 
\pages 1--32



\RBibitem{ruz:bib:5} 
\by Смирнов~М.~М.
\book Уравнения смешанного типа
\publaddr М.
\publ Высшая школа
\yr 1985
\totalpages 304
%Smirnov, M.M.
%Equations of mixed type. (English) Zbl 0377.35001
%Translations of Mathematical Monographs. Vol. 51. Providence, R.I.: American Mathematical Society (AMS). III, 232 p. $ 27.20 (1977).


\RBibitem{ruz:bib:6} 
\by Лернер~М.~Е., Пулькин~С.~П.
\paper О единственности решений задач с условиями Франкля и Трикоми для общего уравнения Лаврентьева"--~Бицадце
\jour Диффер. уравн. 
\yr 1966
\vol 11
\issue 9 
\pages 1255--1263
\zmath{https://zbmath.org/?q=an:0152.09803}


\RBibitem{ruz:bib:7} 
\by Рузиев~М.~Х.
\paper О нелокальной задаче для уравнения смешанного типа с~сингулярным коэффициентом в~неограниченной области
\jour Изв. вузов. Матем.
\yr 2010
\issue 11
\pages 41--49
\mathnet{http://mi.mathnet.ru/ivm7149}
\mathscinet{http://www.ams.org/mathscinet-getitem?mr=2814563}
\zmath{https://zbmath.org/?q=an:1270.35330}
%\transl
%\jour Russian Math. (Iz. VUZ)
%\yr 2010
%\vol 54
%\issue 11
%\pages 36--43
%\crossref{http://dx.doi.org/10.3103/S1066369X10110046}
%Ruziev, M.Kh.
%Author Profile
%A nonlocal problem for a mixed-type equation with a singular coefficient in an unbounded domain. (English. Russian original) Zbl 1270.35330
%Russ. Math. 54, No. 11, 36-43 (2010); translation from Izv. Vyssh. Uchebn. Zaved., Mat. 2010, No. 11, 41-50 (2010).


\RBibitem{ruz:bib:8} 
\by Мирсабуров~М, Рузиев~М.~Х.
\paper Об одной краевой задаче для одного класса уравнений смешанного типа в неограниченной области
\jour Диффер. уравн.
\yr 2011 
\vol 47
\issue 1 
\pages 112--119
\zmath{https://zbmath.org/?q=an:1218.35153}
%Mirsaburov, M.; Ruziev, M.Kh.
%A boundary value problem for a class of mixed-type equations in an unbounded domain. (English. Russian original) Zbl 1218.35153
%Differ. Equ. 47, No. 1, 111-118 (2011); translation from Differ. Uravn. 47, No. 1, 112-119 (2011).


\RBibitem{ruz:bib:9} 
\by Салахитдинов~М.~С., Мирсабуров~М.
\book Нелокальные задачи для уравнений смешанного типа с сингулярными коэффициентами
\publaddr Ташкент
\publ Университет
\yr 2005 
\totalpages 224
%
%Authors of Document Salakhitdinov, M.S., Mirsaburov, M.
%Year the Document was Publish 2005
%Source of the Document Nonlocal Problems for Mixed-Type Equations with Singular Coefficients 


\RBibitem{ruz:bib:10} 
\by Бицадзе~А.~В.
\book Некоторые классы уравнений в частных производных
\publaddr М.
\publ Наука 
\yr 1981
\totalpages 448
\zmath{https://zbmath.org/?q=an:0511.35001}
%Bitsadze, A.V.
%Some classes of partial differential equations. (Nekotorye klassy uravnenij v chastnykh proizvodnykh). (Russian) Zbl 0511.35001
%Moskva: ”Nauka”. 448. p. R. 2.80 (1981).
%Bitsadze, A.V.
%Some classes of partial differential equations. Transl. from the Russian by H. Zahavi. (English) Zbl 0749.35002
%Advanced Studies in Contemporary Mathematics. 4. New York etc.: Gordon & Breach Science Publishers. xi, 504 p. (1988).



\RBibitem{ruz:bib:11} 
\by Мусхелишвили~Н.~И.
\book Сингулярные интегральные уравнения. Граничные задачи теории функций и некоторые их приложения к математической физике
\publaddr М.
\publ Наука 
\yr 1968
\zmath{https://zbmath.org/?q=an:0174.16202}
\totalpages 551
%Muskhelishvili, N.I.
%Some basic problems of the mathematical theory of elasticity. Fundamental equations, plane theory of elasticity, torsion and bending. Translated from the Russian by J. R. M. Radok. 4th, corr. and augm. ed. (English) Zbl 0297.73008
%Leyden: Noordhoff International Publishing. XXXI, 732 p. Dfl. 150.00 (1975).
%10.1007/978-94-017-3034-1


\RBibitem{ruz:bib:12} 
\by Гахов~Ф.~Д., Черский~Ю.~И.
\book Уравнения типа свертки
\publaddr М.
\publ Наука 
\yr 1978
\totalpages 295
\zmath{https://zbmath.org/?q=an:0458.45002}
%Gakhov, F.D.; Cherski?, Yu.I.
%Author Profile
%Equations of convolution type. (Уравнения типа свертки.) (Russian) Zbl 0458.45002
%Moskva: ”Nauka”. 295 p. R. 1.50 (1978).


\end{thebibliography}



\RuDate	

\newpage

\makeentitle

\thanks{{\bf ORCID}

{\bf Menglibay Ruziev:} \url{http://orcid.org/0000-0002-1097-0137}
}



\EngRef




\begin{thebibliography}{99}


\Bibitem{ruz:bib:1-eng} 
\zmath{https://zbmath.org/?q=an:1260.35080}
\by Repin~O.~A., Kumykova~S.~K.
\paper On a boundary value problem with shift for an equation of mixed type in an unbounded domain
\jour Differ. Equ.
\vol 48
\issue  8
\pages 1127--1136 
\yr 2012
\scopus{2-s2.0-84869437089}
\crossref{10.1134/S0012266112080083}

\Bibitem{ruz:bib:2-eng} 
\lang In Russian
\by Abashkin~A.~A.
\paper On one problem in an infinity half-strip for~biaxisimmetric Helmholtz equation
\jour Vestn. Samar. Gos. Tekhn. Univ. Ser. Fiz.-Mat. Nauki
\yr 2012
\pages 39--45
\mathnet{http://mi.mathnet.ru/vsgtu1023}
\crossref{10.14498/vsgtu1023}
\issue 1(26) 



\Bibitem{ruz:bib:3-eng} 
\lang In Russian
\yr 2009 
\issue 1(18) 
\pages 33--40
\by Ruziev~M.~Kh.
\paper Boundary Value Problem For the Equation of Mixed Type with Singular Coefficient in the Domain where the Elliptical Part Is a Half-band
\jour Vestn. Samar. Gos. Tekhn. Univ. Ser. Fiz.-Mat. Nauki
\mathnet{http://mi.mathnet.ru/vsgtu645}
\crossref{10.14498/vsgtu645}



\Bibitem{ruz:bib:4-eng} 
\by Gellerstedt~S. 
\paper Quelques probl\`emes mixtes pour l'\'equation ${y^m}{z_{xx}} + {z_{yy}} = 0$ 
\jour Ark. Mat. Astron. Fys.
\yr 1937
\vol 26A
\issue 3 
\pages 1--32



\Bibitem{ruz:bib:5-eng} 
\by Smirnov~M.~M.
\book Equations of mixed type
\zmath{https://zbmath.org/?q=an:0377.35001}
\serial Translations of Mathematical Monographs
\vol 51
\publaddr Providence, R.I.
\publ  American Mathematical Society (AMS)
\totalpages iii+232 
\yr 1977


\Bibitem{ruz:bib:6-eng} 
\lang In Russian
\by Lerner~M.~E., Pul'kin~S.~P. 
\paper  On uniqueness of solutions of problems satisfying the Frankl and Tricomi conditions for the general Lavrent'ev--Bitsadze equation 
\jour Differ. Uravn. 
\yr 1966
\vol 11
\issue 9 
\pages 1255--1263



\Bibitem{ruz:bib:7-eng} 
\mathscinet{http://www.ams.org/mathscinet-getitem?mr=2814563}
\zmath{https://zbmath.org/?q=an:1270.35330}
\jour Russian Math. (Iz. VUZ)
\yr 2010
\vol 54
\issue 11
\pages 36--43
\crossref{10.3103/S1066369X10110046}
\by Ruziev~M.~Kh.
\paper A nonlocal problem for a mixed-type equation with a singular coefficient in an unbounded domain
\scopus{2-s2.0-78649583445}


\Bibitem{ruz:bib:8-eng} 
\zmath{https://zbmath.org/?q=an:1218.35153}
\by Mirsaburov~M., Ruziev~M.~Kh.
\paper A boundary value problem for a class of mixed-type equations in an unbounded domain
\jour Differ. Equ. 
\vol 47
\issue 1
\pages 111--118 
\yr 2011
\crossref{10.1134/S0012266111010125}
\scopus{2-s2.0-79952380558}



\Bibitem{ruz:bib:9-eng} 
\lang In Russian
\publaddr Tashkent
\publ Universitet
\yr 2005 
\totalpages 224
\by Salakhitdinov~M.~S., Mirsaburov~M.
\book Nelokal'nye zadachi dlia uravnenii smeshannogo tipa s singuliarnymi koeffitsientami \rm [Source of the Document Nonlocal Problems for Mixed-Type Equations with Singular Coefficients]


\Bibitem{ruz:bib:10-eng} 
\zmath{https://zbmath.org/?q=an:0749.35002}
\by Bitsadze~A.~V.
\book Some classes of partial differential equations
\serial Advanced Studies in Con\-tem\-po\-rary Mathematics
\vol 4
\publaddr New York etc.
\publ Gordon \& Breach Science Publishers
\totalpages xi+504 
\yr 1988



\Bibitem{ruz:bib:11-eng} 
\zmath{https://zbmath.org/?q=an:0297.73008}
\by Muskhelishvili~N.~I.
\book Some basic problems of the mathematical theory of elasticity. Fun\-damen\-tal equations, plane theory of elasticity, torsion and bending
\publaddr Leyden
\publ Noordhoff International Publishing
\totalpages xxxi+732 
\yr 1975
\crossref{10.1007/978-94-017-3034-1}


\Bibitem{ruz:bib:12-eng} 
\lang In Russian
\by Gakhov~F.~D., Cherskii~Yu.~I.
\book Uravneniia tipa svertki \rm [Equations of convolution type]
\publaddr Moscow
\publ Nauka 
\yr 1978
\totalpages 295




\end{thebibliography}


\EngDate	


\newpage \Russian

%\end{document}


